\section{Education}

\textbf{University of Electronic Science and Technology of China (UESTC)} \hfill Sept 2022 --- June 2026

\textbf{University of Glasgow, Dual Degree Program} \hfill Sept 2022 --- June 2026

\begin{itemize}
    \item \textbf{Major}: Electrical \& Computer Engineering BEng; \MYhref{https://marcobisky.github.io/cv/score.pdf}{GPA: 3.87/4.0}, \MYhref{https://marcobisky.github.io/cv/rank.pdf}{Ranking: 2/164 (Top 1.2\%)}
    \item \textbf{Relevant Coursework}: Information Theory, Stochastic Processes, Flow Matching and Diffusion Models, Reinforcement Learning in LLM, etc.
\end{itemize}

\textbf{The Chinese University of Hong Kong, Shenzhen} \hfill Sept 2026

\begin{itemize}
    \item \textbf{Major}: Computer Science MPhil
\end{itemize}


\section{Research \& Projects}

\MYhref{https://github.com/Marcobisky/llmlab}{\textbf{LLMlab: Rapid RLVR Post-Training Verification Platform Based on Formal Languages}} \hfill June 2026

\begin{itemize}
    \item Motivation: To reduce the prohibitively high cost of validating LLM training algorithms, LLMlab enables efficient algorithm performance verification on controlled synthetic data.
    \item Designed a deterministic formal language supporting efficient RLVR algorithm benchmarking across varying levels of language difficulty.
    \item Implemented a complete pretrain, SFT, GRPO, KD, OPD, and SDPO pipeline with parallel ablation experiments on 2.67M and 0.15M teacher and student models, achieving 100\% accuracy on difficulty levels 0--3 and over 70\% accuracy on difficulty levels 4--6.
    \item Integrated a visualization toolkit including PCA-projected loss landscapes with weight trajectory tracking, per-layer attention heatmaps, and exposure bias measurements for in-depth model interpretability and training dynamics analysis.
\end{itemize}

\textbf{LLM Post-Training: Reproducing \MYhref{https://arxiv.org/abs/2601.20802}{Self-Distillation Policy Optimization (SDPO)} Paper} \hfill May 2026

\begin{itemize}
    \item Successfully reproduced the SDPO algorithm using the verl RL framework on RunPod H100 GPUs, building on a thorough understanding of the algorithm.
    \item Trained Qwen2.5-3B on code generation tasks using LeetCode-style feedback (runtime errors, failed test cases) as the learning signal, and tracked 40 steps of SDPO training dynamics via WandB, logging mean reward and token-level KL divergence relative to the reference policy.
\end{itemize}

\textbf{System-level Co-Design of RISCV Accelerators for TinyML at the Edge} \hfill Sept 2025 --- April 2026

\vspace{-0.2em} \begin{small} \textit{Research Assistant, \MYhref{https://scholar.google.com/citations?user=1NT1jFMAAAAJ&hl=en}{Prof. Yun Li}, UESTC} \end{small} \vspace{-0.2em}

\begin{itemize}
    \item Engineered a standalone Neural Processing Unit (NPU) for real-time YOLOv8n edge inference on an Artix-7 FPGA, bypassing soft-core processors via a custom RISC-V instruction extension (Xnpu).
    \item Architected an end-to-end Python ML compiler for automated INT16 quantization and memory-aware instruction scheduling, preserving accuracy within 0.3\% mAP of the PyTorch FP32 baseline.
    \item Designed parameterized RTL operators featuring a 3×3 systolic MAC grid and fully hardware-accelerated post-processing (DFL, NMS), achieving 288 MACs/cycle and 23.4 GOPS peak throughput.
    \item Integrated asynchronous camera/UDP video pipelines and AXI4 DDR3L memory multiplexing, fully verified via \texttt{Cocotb}, \texttt{Bazel}, and \texttt{Icarus Verilog}.
\end{itemize}

\MYhref{https://github.com/Marcobisky/my-riscv}{\textbf{Design and Visualization of a Complete Single-cycle RV32I CPU Core}} \hfill Jan 2025 --- Mar 2025

\begin{itemize}
    \item Designed a single-core, single-cycle RISCV 32-bit CPU from scratch in \texttt{Verilog} for RTL simulation and in \texttt{Digital} Software for working principle visualization, open-sourced on \MYhref{https://github.com/Marcobisky/my-riscv}{Github}.
    \item Built a complete datapath including PC, fetcher, decoder, register file, ALU, LRU-based L1 cache, etc., compatible with basic peripherals: GPIOs, IIC, UART, etc.
    \item Implemented a boot program in RISCV assembly, basic delay and GPIO libraries in \texttt{C}. Compiled and simulated using RISCV GNU toolchain.
\end{itemize}


\section{Relevant Skills}

\begin{tabularx}{\linewidth}{@{}l X@{}}
\textbf{Programming} &  \normalsize{\texttt{Python}, \texttt{PyTorch}, \texttt{C/C++} , \texttt{Makefile}.}\\
\textbf{Language} &  \normalsize{Native Chinese, Fluent English (\MYhref{https://marcobisky.github.io/cv/ielts-score.pdf}{IELTS 7}).}\\
\end{tabularx}


\section{Awards}

\textbf{Top Academic Scholarship of UESTC (Top 5\%)} \hfill Dec 2023, Dec 2024 \\
\textbf{China National Scholarship (Top 0.2\%)} \hfill Dec 2024 \\
\textbf{\MYhref{https://www.gla.uestc.edu.cn/info/1089/17553.htm}{Outstanding Graduate of Sichuan Province, 2026}} \hfill Oct 2025

